% TODO make hw stylesheet
\documentclass[10pt]{article}
\usepackage[letterpaper,top=1in]{geometry}
\usepackage[dvipsnames,svgnames]{xcolor}
\usepackage[shortlabels]{enumitem}
\usepackage[framemethod=TikZ]{mdframed}
\usepackage{amsmath,amssymb,amsthm}
\usepackage[colorlinks]{hyperref}
\usepackage{multicol}
\usepackage{tabularx}

\hypersetup{
	colorlinks=true,
	linkcolor=blue,
	filecolor=magenta,
	urlcolor=magenta,
	pdftitle={MAT 528 HW}
}

\usepackage{mathtools}
\usepackage{thmtools}
\usepackage[textsize=footnotesize]{todonotes}
\usepackage{titling}

\reversemarginpar
\mdfdefinestyle{mdbluebox}{roundcorner=10pt,innerbottommargin=9pt,
	linecolor=blue,backgroundcolor=TealBlue!5,}
\declaretheoremstyle[headfont=\sffamily\bfseries\color{MidnightBlue},
mdframed={style=mdbluebox},]{thmbluebox}

\mdfdefinestyle{mdbloobox}{frametitlefont=\bfseries,innerbottommargin=8pt,
	nobreak=true,backgroundcolor=TealBlue!5,linecolor=blue,}
\declaretheoremstyle[headfont=\bfseries\color{MidnightBlue},
mdframed={style=mdbloobox},headpunct={\\[3pt]},postheadspace=0pt,]{thmbloobox}

\mdfdefinestyle{mdredbox}{frametitlefont=\bfseries,innerbottommargin=8pt,
	nobreak=true,backgroundcolor=Salmon!5,linecolor=RawSienna,}
\declaretheoremstyle[headfont=\bfseries\color{RawSienna},
mdframed={style=mdredbox},headpunct={\\[3pt]},postheadspace=0pt,]{thmredbox}

\mdfdefinestyle{mdgreenbox}{linecolor=ForestGreen,backgroundcolor=ForestGreen!5,
	linewidth=2pt,rightline=false,leftline=true,topline=false,bottomline=false,}
\declaretheoremstyle[headfont=\bfseries\sffamily\color{ForestGreen!70!black},
mdframed={style=mdgreenbox},headpunct={ --- },]{thmgreenbox}

\mdfdefinestyle{mdorangebox}{linecolor=RedOrange,backgroundcolor=RedOrange!5,
	linewidth=2pt,rightline=false,leftline=true,topline=false,bottomline=false,}
\declaretheoremstyle[headfont=\bfseries\sffamily\color{RedOrange!70!black},
mdframed={style=mdorangebox},headpunct={ --- },]{thmorangebox}

\mdfdefinestyle{mdblackbox}{linecolor=black,backgroundcolor=RedViolet!5!gray!5,
	linewidth=3pt,rightline=false,leftline=true,topline=false,bottomline=false,}
\declaretheoremstyle[mdframed={style=mdblackbox}]{thmblackbox}

\declaretheoremstyle[spaceabove=6pt,spacebelow=6pt]{basehead}

\declaretheorem[style=basehead,name=Problem]{problem}
\declaretheorem[style=thmredbox,name=Problem,numbered=no]{problem*}
\declaretheorem[style=thmbloobox,name=Setup]{setup} % for config geo
\declaretheorem[style=thmredbox,name=Example,sibling=problem]{example}
\declaretheorem[style=thmredbox,name=$\bigstar$ Problem,sibling=problem]{niceprob}
\declaretheorem[style=thmbluebox,name=Theorem,sibling=problem]{theorem}
\declaretheorem[style=thmbluebox,name=Corollary,sibling=theorem]{corollary}
\declaretheorem[style=thmbluebox,name=Proposition,sibling=theorem]{prop}
\declaretheorem[style=thmbluebox,name=Lemma,numberwithin=problem]{lemma}
\declaretheorem[style=thmbluebox,name=Theorem,numbered=no]{theorem*}
\declaretheorem[style=thmbluebox,name=Lemma,numbered=no]{lemma*}
\declaretheorem[style=thmgreenbox,name=Claim,numberwithin=problem]{claim}
\declaretheorem[style=thmgreenbox,name=Claim,numbered=no]{claim*}
\declaretheorem[style=thmblackbox,name=Remark,numberwithin=problem]{remark}
\declaretheorem[style=thmblackbox,name=Remark,numbered=no]{remark*}
\declaretheorem[style=thmgreenbox,name=Definition,sibling=theorem]{definition}
\declaretheorem[style=thmgreenbox,name=Definition,numbered=no]{definition*}

\providecommand{\ol}{\overline}
\providecommand{\ceil}[1]{\left\lceil #1 \right\rceil}
\providecommand{\floor}[1]{\left\lfloor #1 \right\rfloor}
\newcommand{\dd}{\mathrm{d}}
\providecommand{\ul}{\underline}
\providecommand{\eps}{\varepsilon}
\providecommand{\half}{\frac{1}{2}}
\providecommand{\CC}{\mathbb C}
\providecommand{\FF}{\mathbb F}
\providecommand{\NN}{\mathbb N}
\providecommand{\QQ}{\mathbb Q}
\providecommand{\RR}{\mathbb R}
\providecommand{\ZZ}{\mathbb Z}
\providecommand{\dg}{^\circ}
\providecommand{\ii}{\item}
\providecommand{\setdiff}{\smallsetminus}
\providecommand{\alert}[1]{{\sffamily\textbf{\textcolor{blue}{#1}}}}
\providecommand{\inv}{^{-1}}
\providecommand{\mb}[1]{\mathbf{#1}}

\newenvironment{soln}{\begin{proof}[Solution]}{\end{proof}}

\providecommand{\pow}{\operatorname{Pow}}
\providecommand{\vocab}[1]{{\textbf{\textcolor{ForestGreen}{#1}}}}
\providecommand{\lcm}{\operatorname{lcm}}
\providecommand{\abs}[1]{\left\lvert #1\right\rvert}
\providecommand{\norm}[1]{\left\|#1\right\|}

\usepackage{mathrsfs}

% place title at top right
\setlength{\droptitle}{-8ex}
\pretitle{\begin{flushright}\Large\bfseries}
\posttitle{\par\end{flushright}}
\preauthor{\begin{flushright}}
\postauthor{\end{flushright}}
\predate{\begin{flushright}}
\postdate{\end{flushright}}

\title{MAT 528 HW09}
\author{\large Aditya Pahuja}
\date{\large Due 11/13}
\begin{document}
\maketitle

\begin{problem}
    [Munkres 33.3]
    Give a direct proof of the Urysohn lemma
    for a metric space $(X,d)$ by setting
    \begin{equation*}
	f(x) = \frac{d(x,A)}{d(x,A) + d(x,B)}.
    \end{equation*}
\end{problem}

\begin{soln}
    Let $A$ and $B$ be disjoint closed sets and $f$ as given.
    Recall that if $C$ is a closed set, then $d(x,C) = 0$
    is equivalent to $x\in C$, for $d(x,C) = 0$ implies
    $B_\eps(x)\cap C$ is nonempty for all $\eps>0$,
    hence $x$ is a limit point of $C$, so belongs to $C$
    (and the converse is trivial).

    First, we show that $f$ is continuous.
    Now observe that for any $C\subseteq X$ and $x,y\in X$,
    \begin{equation*}
	d(x,C) = \inf_{z\in C} d(x,z) \le \inf_{z\in C}\{d(x,y) + d(y,z)\}
	= d(x,y) + d(y,C)
    \end{equation*}
    and symmetrically $d(y,C) \le d(x,y) + d(x,C)$,
    so $\abs{d(x,C) - d(y,C)} \le d(x,y)$.
    Thus, if $d(x,y) < \eps$ then so is $\abs{d(x,C) - d(y,C)}$,
    proving that $d(x,C)$ is continuous in $x$,
    meaning $f$ is continuous at all $x\in X$ such that $d(x,A) + d(x,B)$ is nonzero.
    Since $d(x,A)$ and $d(x,B)$ are nonnegative,
    this can happen only if $d(x,A) = d(x,B) = 0$,
    i.e. $x\in A\cap B = \varnothing$, which is to say that
    $d(x,A) + d(x,B)$ is never zero. Therefore $f$ is continuous on $X$.

    To see that $f$ separates $A$ and $B$, we see that
    $x\in A$ implies $d(x,A) = 0$, so
    \begin{equation*}
	f(x) = \frac{0}{0 + d(x,B)} = 0,
    \end{equation*}
    while $x\in B$ implies that $d(x,B) = 0$ and since $x\notin A$,
    $d(x,A)\ne 0$, hence
    \begin{equation*}
	f(x) = \frac{d(x,A)}{d(x,A) + 0} = 1.\qedhere
    \end{equation*}
\end{soln}

\begin{problem}
    [Munkres 36.1]
    Prove that every manifold is regular and hence metrizable.
    Where do you use the Hausdorff condition?
\end{problem}

\begin{soln}
    Let $M$ be an $m$-manifold, $x\in M$, and $U\subseteq M$
    a neighborhood of $x$. Then, there exists an open
    $W\subseteq M$ that is homeomorphic to open $V\subseteq\RR^m$
    via homeomorphism $h$; take $W' = W\cap U$
    and $V' = h(W')$ the homeomorphic image of $W'$.
    Since $V'$ is open it contains a closed ball $C$ centered at $h(x)$
    (which, by Heine-Borel, is compact),
    and $C$ contains an open ball $B$ centered at $h(x)$.
    Since $h\inv|_{V'}\colon V'\to W'$ is continuous,
    $h\inv(C)\subseteq W'$ and
    $h\inv(B)\subseteq h\inv(C)$ are respectively compact and open,
    both containing $x$.

    I claim that the closure of $h\inv(B)$
    is contained in $U$: since $h\inv(C)$ is compact, it is closed in $M$,
    and by construction it contains $h\inv(B)$,
    so $h\inv(C)$ contains the minimal closed set (under inclusion)
    containing $h\inv(B)$, which by definition is the closure of $h\inv(B)$.

    To summarize, we have shown that, given $x\in M$
    and a neighborhood $U\subseteq M$ of $x$,
    there exists a neighborhood of $x$ whose closure lies in $U$,
    which by Munkres's Lemma 31.1 is equivalent to $M$ being a regular space.
\end{soln}

\begin{problem}
    [Munkres 43.5]
    Show that if $f$ is a contraction of a complete metric space,
    then there is a unique point $x\in X$ such that $f(x) = x$.
\end{problem}

\begin{soln}
    First, we show that $f$ has a fixed point.
    Fix $x_0\in X$ and let $x_n = f(x_{n-1})$ for positive integer $n$.
    Since $f$ is a contraction, there exists $\alpha \in (0,1)$ such that
    $d(f(x), f(y)) \le \alpha d(x,y)$ for all $x,y\in X$.
    If $f^k$ is the function $f$ composed with itself $k$ times, then
    induction tells us $d(f^k(x), f^k(y)) \le \alpha^k d(x,y)$.
    Letting $x_0\in X$ be arbitrary and setting $x_n = f^n(x_0)$ for $n\in\NN$,
    we get that for integers $m < n$,
    \begin{equation*}
	d(x_n, x_m) \le \sum_{j=m}^{n-1} d(x_j, x_{j+1})
	\le \sum_{j=m}^{n-1} \alpha^j d(x_0,x_1)
	\le d(x_0,x_1)\cdot \sum_{j=m}^\infty \alpha^j
	= d(x_0,x_1)\cdot \frac{\alpha^m}{1-\alpha}.
    \end{equation*}
    Since $\lim_{m\to\infty} \alpha^m = 0$,
    given $\eps > 0$ there exists $N$ such that
    $d(x_0,x_1)\cdot \alpha^N/(1-\alpha) < \eps$, so that
    $n > m > N$ implies
    \begin{equation*}
        d(x_n,x_m)
	\le d(x_0,x_1)\cdot \frac{\alpha^m}{1-\alpha}
	\le d(x_0,x_1)\cdot \frac{\alpha^N}{1-\alpha}
	\le \eps.
    \end{equation*}
    This means $(x_n)$ is a Cauchy sequence,
    so it converges to some $x_\infty\in X$.
    This $x_\infty$ is a fixed point of $f$, since
    \begin{equation*}
	x_\infty = \lim_{n\to\infty} x_n
	= \lim_{n\to\infty}f(x_{n-1})
	= f\left(\lim_{n\to\infty} x_n\right)
	= f(x_\infty)
    \end{equation*}
    where the limit can be pulled inside $f$ because
    contractions are continuous (where $d(x,y) < \eps$
    implies $d(f(x), f(y)) \le d(x,y) < \eps$).

    The fixed point is unique because
    if $x,x'\in X$ are distinct fixed points of $f$, then
    \begin{equation*}
        d(x,x') \le \alpha d(f(x),f(x')) = \alpha d(x,x') < d(x,x'),
    \end{equation*}
    which is impossible.
\end{soln}

\begin{problem}
    [Munkres 43.7]
    Show that the set of all sequences $(x_1,x_2,\dots)$
    such that $\sum_{k=1}^\infty x_k^2$ converges is complete in the $\ell^2$ metric.
\end{problem}

\begin{soln}
    The set of all such sequences is a normed vector space
    with norm $\norm{\mb x} = (\sum_{k=1}^\infty x_k^2)^{1/2}$,
    so $d(\mb x,\mb y) = \norm{\mb x-\mb y}$ (as proven in HW03).
    Let $\mb x_n = (x_{n,1}, x_{n,2}, \dots)$, $n\in\NN$,
    be a Cauchy sequence in $\ell^2$.
    Each of the sequences $(x_{m,k})_{m\in\NN}$ is Cauchy
    because $\abs{x_{m,k} - x_{n,k}} < \norm{\mb x_m - \mb x_n} < \eps$
    for large enough $m$, $n$, so $\mb x_n$ converges coordinatewise
    to some $\mb x = (\ell_1,\ell_2,\dots)$
    (i.e. $\ell_k = \lim_{m\to\infty} x_{m,k}$).

    % For each $r\in\NN$ and $\eps>0$, there exists $N_{r,\eps}$ such that
    % $n > N_{r,\eps}$ implies $\sum_{k=1}^r \abs{x_k - x_{n,k}}^2 < \eps$
    % because the $x_{n,k}$ converge to $x_k$.
    % In the limit, $\sum_{k=1}^\infty |x_k - x_{n,k}^2| \le \eps$.

    \begin{claim}
	The vector
	\[\mb x = (\ell_1, \ell_2, \ell_3, \dots)\]
	is in $\ell^2$.
    \end{claim}
    \begin{proof}
	We wish to show that
	\[\sum_{i=1}^\infty |\ell_i|^2\]
	converges. It suffices to show that
	the partial sums are bounded, because the numbers being summed
	are all nonnegative.
	For each $i\le N$, there exists an $M_i$ such that
	\[n > M_i\implies |\ell_i - x_{n,i}| < 2^{-i}.\]
	Fix $n > \max\{M_1, \dots, M_N\}$.
	The $N$th partial sum can be bounded above as follows:
	\begin{equation*}
	    \left(\sum_{i=1}^N |\ell_i|\right)^{1/2}
	    \le \left(\sum_{i=1}^N |\ell_i - x_{n,i}|\right)^{1/2}
	    + \left(\sum_{i=1}^N |x_{n,i}|\right)^{1/2}
	    < \left(\sum_{i=1}^\infty 2^{-i}\right)^{1/2} + \|\mb x_n\|
	    = 1 + \|\mb x_n\|.
	\end{equation*}
	Additionally, Cauchy sequences are bounded,
	so $\|\mb x_n\|$ is bounded above by a constant $C$
	for all $n$. Therefore we have bounded the partial sums of the $|\ell_i|^2$
	by $(1 + C)^2$, so the series converges.
    \end{proof}

    Now, we need to prove that $\mb x_n\to\mb x$
    with respect to the $\ell^2$ metric. Let $\eps>0$ and $N\in\NN$.
    For each $i=1,\dots,N$, we can find thresholds $M_i$ for which
    \[m > M_i \implies |\ell_i - x_{m,i}|^2 < \eps^2\cdot 2^{-i - 2}.\]
    This means that, if $m > \max\{M_1, \dots, M_N\}$, then
    \[\left(\sum_{i=1}^N |\ell_i - x_{m,i}|^2\right)^{1/2}
    < \left(\sum_{i=1}^\infty \eps^2\cdot 2^{-i-2}\right)^{1/2}
    = \sqrt{\eps^2/4} = \eps/2\]
    We also know that $(\mb x_n)$ being Cauchy means that
    there is an $M$ such that
    $m,n > M$ implies $\|\mb x_m - \mb x_n\| < \frac\eps2$.
    Putting this together, we can fix
    $m$ larger than $M$ and all of the $M_i$
    and use the triangle inequality to say that
    $n > M$ implies
    \begin{align*}
	\left(\sum_{i=1}^N |\ell_i - x_{n,i}|\right)^{1/2}
	&\le \left(\sum_{i=1}^N |\ell_i - x_{m,i}|\right)^{1/2} +
	\left(\sum_{i=1}^N |x_{m,i} - x_{n,i}|\right)^{1/2}\\
	&\le \eps/2 + \norm{x_m - x_n} < \eps.
    \end{align*}
    In other words, as long as $n > M$ (noting that $M$ only depends on $\eps$),
    we can bound all the partial sums of $N$ below $\eps$,
    so $n > M$ implies that $\|\mb x - \mb x_n\| < \eps$,
    which is what it means for $\mb x_n$ to converge to $\mb x$
    in the metric of $\ell^1$, so we are done.
\end{soln}

\begin{problem}
    [Munkres 43.9]
    Let $(X,d)$ be a metric space. Let $\tilde X$ denote the set
    of all Cauchy sequences $\mb x = (x_1,x_2,\dots)$ of points of $X$.
    Define $\mb x\sim\mb y$ if $d(x_n,y_n)\to 0$.
    Let $[\mb x]$ denote the equivalence class of $\mb x$,
    and let $Y$ denote the set of these equivalence classes.
    Define a metric $D$ on $Y$ by the equation
    $D([\mb x],[\mb y]) = \lim_{n\to\infty} d(x_n,y_n)$.
    \begin{enumerate}[(a)]
	\ii Show that $\sim$ is an equivalence relation,
	and show that $D$ is a well-defined metric.
	\ii Define $h\colon X\to Y$ by letting $h(x)$ be
	the equivalence class of the constant sequence $(x,x,\dots)$.
	Show that $h$ is an isometric embedding.
	\ii Show that $h(X)$ is dense in $Y$;
	indeed, given $\mb x = (x_1,x_2,\dots)\in\tilde X$,
	show that the sequence $h(x_n)$ of points of $Y$
	converges to the point $[\mb x]$ of $Y$.
	\ii Show that if $A$ is a dense subset of a metric space $(Z,\rho)$
	and if every Cauchy sequence in $A$ converges in $Z$,
	then $Z$ is complete.
	\ii Show that $(Y,D)$ is complete.
\end{enumerate}
\end{problem}

\begin{soln}
    (a) Let $\mb x,\mb y,\mb z\in\tilde X$.
    The relation is reflexive, since $\mb x\sim\mb x$ because
    $\lim_{n\to\infty} d(x_n,x_n) = \lim_{n\to\infty} 0 = 0$.
    The relation is symmetric because
    $\lim_{n\to\infty}d(x_n,y_n) = 0$ implies $\lim_{n\to\infty} d(y_n,x_n) = 0$
    hence $\mb y\sim\mb x$.
    The relation is transitive because if $\mb x\sim\mb y$
    and $\mb y\sim\mb z$ then
    \begin{equation*}
	\lim_{n\to\infty} d(x_n,z_n)
	\le \lim_{n\to\infty} d(x_n,y_n) + d(y_n,z_n)
	= \lim_{n\to\infty}d(x_n,y_n) + \lim_{n\to\infty}d(y_n,z_n) = 0,
    \end{equation*}
    so the limit on the left is zero, i.e. $\mb x\sim\mb z$.

    To check that $D$ is well-defined we need to check that
    if $\mb x_1\sim\mb x_2$ and $\mb y_1\sim\mb y_2$ then
    \begin{equation*}
	\lim_{n\to\infty}d(x_{1,n},y_{1,n})
	= \lim_{n\to\infty}d(x_{2,n},y_{2,n})
    \end{equation*}
    where $\mb x_j = (x_{j,1}, x_{j,2}, \dots)$
    and $\mb y_j = (y_{j,1}, y_{j,2}, \dots)$ for $j=1,2$.
    By the triangle inequality,
    \begin{align*}
	d(x_{1,n},y_{1,n}) &\le
	d(x_{1,n},x_{2,n}) + d(x_{2n},y_{2,n}) + d(y_{2,n},y_{1,n})\\
	d(x_{2,n},y_{2,n}) &\le
	d(x_{2,n},x_{1,n}) + d(x_{1n},y_{1,n}) + d(y_{1,n},y_{2,n}),
    \end{align*}
    so we have the upper bound
    \begin{equation*}
	\abs{d(x_{1,n},y_{1,n}) - d(x_{2,n},y_{2,n})}
	\le d(x_{1,n},x_{2,n}) + d(y_{1,n},y_{2,n}).
    \end{equation*}
    The right-hand side converges to zero because $\mb x_1\sim\mb x_2$
    and $\mb y_1\sim\mb y_2$, hence so does the left-hand side,
    meaning $d(x_{1,n},y_{1,n})$ and $d(x_{2,n},y_{2,n})$ have the same limit.
    \\

    (b) For any $x,y\in X$,
    $D(h(x), h(y)) = \lim_{n\to\infty} d(x,y) = d(x,y)$,
    so $h$ is an isometric embedding.
    \\

    (c) Since $\mb x$ is a Cauchy sequence,
    given $\eps>0$, there exists $N_\eps\in\NN$ such that $m,n>N_\eps$ implies
    $d(x_m,x_n) < \eps$. Then, if $n > N_\eps$,
    \begin{equation*}
	D(h(x_n), \mb x) = \lim_{k\to\infty} d(x_n,x_k) < \eps
    \end{equation*}
    because the argument of the limit will be smaller than $\eps$
    for all $k > N_\eps$, which means $h(x_n)$ converges to $\mb x$.
    This implies $\mb x$ is a limit point of $h(X)$,
    so each point in $Y$ is a limit point of $h(X)$,
    to elaborate on why $h(X)$ is dense in $Y$.
    \\

    (d) Let $(x_n)$ be a Cauchy sequence in $Z$
    and, for $\eps>0$, let $N_\eps\in\NN$ be such that
    $m,n>N_\eps$ implies $\rho(x_m,x_n) < \eps$.
    Since $A$ is dense in $Z$, for each $n\in\NN$, there exists
    $a_n\in A$ such that $\rho(a_n,x_n) < 1/n$
    (since the open ball $B_{1/n}(x_n)$ must intersect $A$);
    since $\lim_{n\to\infty} \rho(a_n,x_n) = 0$,
    if either sequence $(a_n)$ or $(x_n)$ converges, then so does the other one
    and both sequences have the same limit.
    The sequence $(a_n)$ is a Cauchy sequence:
    if $\eps>0$, pick $N\in\NN$ such that $N > N_{\eps/3}$ and
    $1/N < \eps/3$, so $m,n>N$ implies
    \begin{equation*}
	\rho(a_m,a_n) \le \rho(a_m,x_m) + \rho(x_m,x_n) + \rho(x_n,a_n)
	< \frac 1m + \frac 1n + \frac\eps3 \le \frac 2N + \frac\eps3
	< \eps.
    \end{equation*}
    By assumption, then, $(a_n)$ converges in $Z$,
    hence so does $(x_n)$, with $\lim_{n\to\infty} a_n = \lim_{n\to\infty} x_n$.
    Thus every Cauchy sequence $(x_n)$ in $Z$
    has a limit in $Z$, so $Z$ is complete.
    \\

    (e) By (c), $h(X)$ is dense in $Y$
    and every Cauchy sequence $(h(x_n))$ in $h(X)$ converges to $[(x_n)]\in Y$,
    so, by (d) with $A = h(X)$ and $Z = Y$, $Y$ is complete.
\end{soln}

\begin{problem}
    [Munkres 45.4a]
    Let $f_n\colon[0,1]\to\RR$ be the function $f_n(x) = x^n$.
    The collection $\mathcal F = \{f_n\}$ is pointwise bounded
    but the sequence $(f_n)$ has no uniformly convergent subsequence;
    at what point or points does $\mathcal F$ fail to be equicontinuous?
\end{problem}

\begin{soln}
    $\mathcal F$ is equicontinuous at all $x_0\in[0,1)$
    and not equicontinuous at $x_0 = 1$.

    To show that $\mathcal F$ is (pointwise) equicontinuous on $[0,1)$,
    we prove (pointwise) equicontinuity over $[0,c]$ for every $c\in[0,1)$.
    Fix $c\in[0,1)$. For each $\eps>0$, there exists $N\in\NN$ such that
    $c^N < \eps$, hence $\abs{x^n - 0} < \eps$ for all $n > N$
    and $x\in[0,c]$; this means that $x^n$ converges to the zero function
    uniformly over $[0,c]$. The zero function is the only nontrivial limit point
    $\mathcal F|_{[0,c]} \coloneq\{f_n|_{[0,c]} : n\in\NN\}$, since
    any limit point would be a limit of a subsequence of the $f_n|_{[0,c]}$,
    which is equal to the limit of the $f_n|_{[0,c]}$ (i.e. the zero function).
    That is to say, the closure of $\mathcal F|_{[0,c]}$
    in the space of continuous functions over $[0,c]$ consists of
    $\mathcal F|_{[0,c]}$ together with the zero function.

    This closure is easily seen to be compact,
    as, given an open cover of $\ol{\mathcal F}|_{[0,c]}$,
    whichever open set contains the zero function must also contain
    all but finitely many elements of the closure;
    and then the other elements can each have their own open set reserved for them
    to obtain a finite subcover.
    Now, Ascoli's theorem (Munkres's Theorem 45.4) tells us directly
    that $\mathcal F|_{[0,c]}$ is equicontinuous, meaning
    $\mathcal F$ is equicontinuous on $[0,c)$ for every $c\in[0,1)$,
    hence equicontinuous on $[0,1)$.

    Now, to show that $\mathcal F$ is not equicontinuous at $1$,
    suppose for contradiction that there exists $U$ open in $[0,1]$ and
    containing $1$ such that $1-x^n < \half$
    holds for all $x\in U$ and $n\in\NN$.
    However, if one picks $x\in U$ not equal to (hence smaller than) $1$,
    then $\lim_{n\to\infty}x^n = 0$, so for each $x\in U\setminus\{1\}$
    there exists $n$ such that $x^n < \half$, i.e. $1 - x^n > \half$;
    this proves that $\mathcal F$ is not equicontinuous at $1$.
\end{soln}










\end{document}

